\documentclass[final]{beamer}  
\mode<presentation> { \usetheme{ParisSaclay} }
\usepackage[english, spanish]{babel}
\usepackage[utf8]{inputenc}
\usepackage{amsmath,amsthm, amssymb, latexsym}

\boldmath %math en gras
\usepackage[orientation=portrait,size=a1,scale=1.4]{beamerposter}

\usepackage{graphicx}
\usepackage{tikz}

% --- Título y Autores Reales ---
\title{\huge INTEGRACIÓN DE SENTIMIENTO DE MERCADO EN AGENTES DE RL PARA LA OPTIMIZACIÓN DE TRADING}
\author{Julian Duarte, Julian Jimenez}
\institute[]{Universidad Externado de Colombia \\ Ciencia de Datos -- Departamento de Matemáticas \\ Semillero de Aprendizaje por Refuerzo \\ \medskip \small Tutores: Prof. Julián Fernando Sánchez \quad Prof. Juan Carlos Urueña}

\begin{document}

\begin{frame} [fragile]{} 

\begin{columns}
    \begin{column}{.48\textwidth}
        
    \begin{block}{1. Introducción y Contexto}
        El mercado de criptomonedas se caracteriza por una volatilidad extrema impulsada por sesgos cognitivos. La toma de decisiones suele verse afectada por estados emocionales como el pánico o la euforia (FOMO).
        \begin{itemize}
            \item \textbf{Problema:} Los modelos tradicionales ignoran el sentimiento colectivo.
            \item \textbf{Necesidad:} Sistemas racionales que integren indicadores emocionales.
        \end{itemize}
    \end{block}

    \vspace{1cm}

    \begin{block}{2. Objetivo(s) del Trabajo}
        Desarrollar y validar un agente de \textbf{Q-Learning} que ejecute estrategias rentables minimizando sesgos emocionales mediante la incorporación del \textit{Fear \& Greed Index}.
    \end{block}

    \vspace{1cm}

    \begin{alertblock}{3. Metodología (MDP)}
        Optimización de una política $\pi$ para maximizar la recompensa acumulada:
        \begin{description}
            \item[Estado ($S$):] RSI, Medias Móviles, Sentimiento e Inventario.
            \item[Acciones ($A$):] $\{ \text{Compra, Venta, Neutro} \}$.
            \item[Recompensa ($R$):] $R_t = \Delta \text{P\&L} - \text{Comisiones} - \lambda \sigma_t$.
        \end{description}
        
        \centering
        \vspace{0.5cm}
        \[ Q(s, a) \leftarrow Q(s, a) + \alpha [r + \gamma \max Q(s', a') - Q(s, a)] \]
        \vspace{0.5cm}
        
        \textbf{[Espacio para Diagrama del Ciclo de RL]}
    \end{alertblock}
    
    \end{column}

    \begin{column}{.48\textwidth}	

    \begin{block}{4. Resultados Esperados y Evaluación}
        Dado que el proyecto se encuentra en fase de validación técnica:
        \begin{itemize}
            \item Se proyecta una reducción de la varianza en la toma de decisiones.
            \item Mayor resiliencia del capital durante periodos de pánico.
            \item \textbf{Métricas:} Ratio de Sharpe, Max Drawdown y Profit Factor.
        \end{itemize}
        
        \centering
        \vspace{1cm}
        \includegraphics[width=.8\textwidth]{example-image} 	
        \captionof{figure}{Simulación Hipotética de Rendimiento}
    \end{block}
    
    \vspace{1cm}	

    \begin{exampleblock}{5. Conclusiones y Aplicación}
        La integración de finanzas conductuales en agentes autónomos permite una gestión de riesgo más robusta, sentando bases para sistemas de ahorro inteligente resilientes a crisis financieras.
    \end{exampleblock}

    \vspace{1cm}

    \begin{block}{6. Referencias y Contacto}
        \begin{enumerate}
            \item Sutton \& Barto, \textit{Reinforcement Learning: An Introduction}, MIT Press.
            \item Alternative.me, \textit{Crypto Fear \& Greed Index API}.
        \end{enumerate}
        \centering
        \vspace{0.5cm}
        \textbf{Contacto:} julian.duarte@uexternado.edu.co
    \end{block}

    \end{column}
\end{columns}

\end{frame}

\end{document}
